(i). On the other hand, (iii) implies (ii) by 2.3 and 2.8.

\sgaStableTarget{sga6.e13.sec.4}\subsection*{4. Finiteness theorems for @@M0@@}

\sgaStableTarget{sga6.e13.para.4-1}\textbf{4.1.} Let @@M1@@ be a scheme and @@M2@@ a commutative group functor on @@M3@@. We denote by @@M4@@ (resp. @@M5@@) the group subfunctor whose @@M6@@-points @@M7@@ are those of @@M8@@ such that, for every algebraically closed field @@M9@@ and every @@M10@@-point @@M11@@ of @@M12@@, there exist a connected @@M13@@-scheme @@M14@@, two @@M15@@-points @@M16@@ of @@M17@@ and a morphism @@M18@@ such that @@M19@@ and @@M20@@ in @@M21@@ (resp. there exists a nonzero integer @@M22@@ such that @@M23@@).

It is evident that the formation of @@M24@@ (resp. @@M25@@) commutes with every base change, that a morphism of group functors @@M26@@ induces a morphism @@M27@@ (resp. @@M28@@), and that if @@M29@@ is a field and @@M30@@ is represented by a group scheme locally of finite type, then @@M31@@ is represented by the connected component of the identity element, which is an open (and closed) subscheme, of finite type and geometrically irreducible (SGA 3 VI@@M32@@ 2.4).

\sgaStableTarget{sga6.e13.lem.4-2}\textbf{Lemma 4.2.} \emph{Let @@M33@@ be a morphism between two commutative group functors on a scheme @@M34@@. We suppose that, for every algebraically closed field @@M35@@ and every morphism @@M36@@, @@M37@@ (resp. @@M38@@) is representable by a group scheme locally of finite type, and @@M39@@ is representable by a morphism of finite type, (that is, quasi-compact). Then}
@@M40@@

